%%--------------------------------------------------------------------
%1--------------------------------------------------------------------
\begin{glyph}{Law (法)}{law}\label{glyph:law}
Law.  Method.\\\hline
A religious (particularly Buddhist) meaning is also conveyed in some compounds that include this 漢字.
\end{glyph}
\gindex{Law}{法}\strindex{8}{法}

\begin{comps}
\comprtwocone&水氵氺 + 去\\\hline
\comprthreecone&Water + Depart (p. \pageref{glyph:depart})\\\hline
\end{comps}

\begin{remglyph}
\hooglig{Law} stipulates that \remcom{water} must \remcom{depart} as urine.\\\hline
\end{remglyph}

\begin{prons}
\pronsrtwocone&\textbf{ホウ (H\={O})}&---\\\hline
\pronsrthreecone&ハッ、ホッ (HA-, HO-)&---\\\hline
\rye{3}{These readings always `double up' with the 音読み of the 漢字 that follow.}
\end{prons}

\begin{remprons}
\hooglig{Law} dictates that \ucomon{harbingers} should \comon{haul} the corpses to \ucomon{hospital}.\\\hline
\end{remprons}

%{\middel}
% &  \mbox{()}\\\hline
\begin{somme}
方法&direction + law = \textit{method; way}&ホウ{\middel}ホウ \mbox{(H\={O}{\middel}H\={O})}\\\hline
文法&literature + law = \textit{grammar}&ブン{\middel}ポウ \mbox{(BUN{\middel}P\={O})}\\\hline
立法&stand + law = \textit{legislation; law making}&リッ{\middel}ポウ \mbox{(RIP{\middel}P\={O})}\\\hline
手法&hand + law = \textit{technique; method}&シュ{\middel}ホウ \mbox{(SHU{\middel}H\={O})}\\\hline
非合法&un- + match + law = \textit{unlawful; illegal}&ヒ{\middel}ゴウ{\middel}ホウ \mbox{(HI{\middel}G\={O}{\middel}H\={O})}\\\hline
法令&law + order = \textit{laws and ordinances}&ホウ{\middel}レイ \mbox{(H\={O}{\middel}REI)}\\\hline
法王&law + king = \textit{Buddha}&ホウ{\middel}オウ \mbox{(H\={O}{\middel}\={O})}\\\hline
説法&opinion + law = \textit{lecture; sermon; preaching}&セッ{\middel}ポウ \mbox{(SEP{\middel}P\={O})}\\\hline
書法&write + law = \textit{penmanship; calligraphy}&ショ{\middel}ホウ \mbox{(SHO{\middel}H\={O})}\\\hline
無作法&nothingness + manners = \textit{ill-mannered; rude}&ブ{\middel}サ{\middel}ホウ \mbox{(BU{\middel}SA{\middel}H\={O})}\\\hline
法要&law + essential = \textit{Buddhist Memorial Service}&ホウ{\middel}ヨウ \mbox{(H\={O}{\middel}Y\={O})}\\\hline
法事&law + matter = \textit{Buddhist Memorial Service}&ホウ{\middel}ジ \mbox{(H\={O}{\middel}JI)}\\\hline
\notyet{泳法}&swim + law = \textit{swimming style}&エイ{\middel}ホウ \mbox{(EI{\middel}H\={O})}\\\hline
\notyet{礼法}&courtesy + law = \textit{etiquette; manners}&レイ{\middel}ホウ \mbox{(REI{\middel}H\={O})}\\\hline
\notyet{自然法}&nature + law = \textit{natural law}&シ{\middel}ゼン{\middel}ホウ \mbox{(SHI{\middel}ZEN{\middel}H\={O})}\\\hline
\end{somme}
\end{multicols}

\begin{decomp}{5}
別&の&方法&を&考え出そう。\\\hline
ベツ&の&ホウホウ&を&かんがえだそう\\\hline
different&の&method&を&to think out\\\hline
\multicolumn{5}{|r|}{\textit{Let's think out another way.}}\\\hline
\end{decomp}
\pagebreak
%%--------------------------------------------------------------------
%2--------------------------------------------------------------------
\begin{multicols}{2}
\begin{glyph}{Place (場)}{place}\label{glyph:place}
A place.
\end{glyph}
\gindex{Place}{場}\strindex{12}{場}

\begin{comps}
\comprtwocone&土 + 昜\\\hline
\comprthreecone&Earth/Soil + Glorious/Bright\\\hline
\rye{2}{昜 = 旦 + 勿.  Glorious = Dawn + {\ldots} Not.  \hooglig{The glorious} \remcom{dawn} of time \remcom{must not} be forgotten.}
\end{comps}

\begin{remglyph}
\hooglig{A place} on \remcom{earth} is not that \remcom{glorious} as a \hooglig{place} in heaven.\\\hline
\end{remglyph}

\begin{prons}
\pronsrtwocone & \textbf{ジョウ (J\={O})} & \textbf{ば (ba)}\\\hline
\pronsrthreecone & --- & --- \\\hline
\end{prons}

\begin{remprons}
\hooglig{Places} \comkun{bulging} with \comon{jaw-dropping} experiences.\\\hline
\end{remprons}

%{\middel}
% &  \mbox{()}\\\hline
\begin{somme}
工場 & craft + place = \textit{factory; plant} & コウ{\middel}ジョウ \mbox{(K\={O}{\middel}J\={O})}\\\hline
場所 & place + location = \textit{place; location} & ば{\middel}ショ \mbox{(ba{\middel}SHO)}\\\hline
場合 & place + match = \textit{circumstances; situation; case} & ば{\middel}あい \mbox{(ba{\middel}ai)}\\\hline
売り場&sell + place = \textit{selling area; sales floor; counter}&う{\middel}り{\middel}ば \mbox{(u{\middel}ri{\middel}ba)}\\\hline
立場&stand + place = \textit{position; location}&たち{\middel}ば \mbox{(tachi{\middel}ba)}\\\hline
道場&road + place = \textit{martial arts training hall}&ドウ{\middel}ジョウ \mbox{(D\={O}{\middel}J\={O})}\\\hline
式場&type + place = \textit{place of ceremony}&シキ{\middel}ジョウ \mbox{(SHIKI{\middel}J\={O})}\\\hline
来場&come + place = \textit{attendance}&ライ{\middel}ジョウ \mbox{(RAI{\middel}J\={O})}\\\hline
\end{somme}

\begin{decomp}{6}
雨&の&場合&私&は&行かない。\\\hline
あめ&の&ばあい&わたし&は&いかない\\\hline
rain&の&circumstance&I&は&not go\\\hline
\multicolumn{6}{|r|}{\textit{In case it rains, I won't go.}}\\\hline
\end{decomp}
\end{multicols}

\begin{multicols}{2}
\begin{glyph}{Desire (願)}{desire}\label{glyph:desire}
Desire.  Petition.\\\hline
This character carries a light tone of entreaty.
\end{glyph}
\gindex{Desire}{願}\strindex{19}{願}

\begin{comps}
\comprtwocone&原 + 頁\\\hline
\comprthreecone&Original (p. \pageref{glyph:original}) + Big shell/Page/Face\\\hline
\end{comps}

\begin{remglyph}
\hooglig{Desire} for the \remcom{original} \remcom{big shells} on this page.\\\hline
\end{remglyph}

\begin{prons}
\pronsrtwocone&\textbf{ガン (GAN)}&\textbf{ねが (nega)}\\\hline
\pronsrthreecone&---&---\\\hline
\end{prons}

\begin{remprons}
\hooglig{Desire} to \comkun{negate} the \comon{Guns}.\\\hline
\end{remprons}

%{\middel}
% &  \mbox{()}\\\hline
\begin{somme}
願う & \textit{desire; request politely}&ねが{\middel}う \mbox{(nega{\middel}u)}\\\hline
お願いします&\textit{please}&お{\middel}ねが{\middel}い{\middel}します \mbox{(o{\middel}nega{\middel}i{\middel}shimasu)}\\\hline
出願者&application + individual = \textit{applicant}&シュツ{\middel}ガン{\middel}シャ (SHUTSU{\middel}GAN{\middel}SHA)\\\hline
宿願&lodging + desire = \textit{longstanding desire}&シュク{\middel}ガン \mbox{(SHUKU{\middel}GAN)}\\\hline
願い下げ&desire + lower = \textit{cancellation; withdrawal}&ねが{\middel}い{\middel}さげ \mbox{(nega{\middel}i{\middel}sa{\middel}ge)}\\\hline
切願&cut + desire = \textit{supplication}&セツ{\middel}ガン \mbox{(SETSU{\middel}GAN)}\\\hline
特願&special + desire = \textit{patent application}&トク{\middel}ガン \mbox{(TOKU{\middel}GAN)}\\\hline
結婚願望{\notcov}&marriage + desire + ambition = \textit{yearning for marriage}&ケッ{\middel}コン{\middel}ガン{\middel}ボウ (KEK{\middel}KON{\middel}GAN{\middel}B\={O})\\\hline
\end{somme}

\begin{decomp}{5}
私達&は&平和&を&願う。\\\hline
わたしたち&は&ヘイワ&を&ねがう\\\hline
we/us&は&peace&を&to desire\\\hline
\multicolumn{5}{|r|}{\textit{We long for peace.}}\\\hline
\end{decomp}

\end{multicols}
\pagebreak
%%--------------------------------------------------------------------
%3--------------------------------------------------------------------
\begin{multicols}{2}
\begin{glyph}{Course (科)}{course}\label{glyph:course}
\textit{Course of study}.  Ideas related to \textit{scholarly matters}.\\\hline
This character is used to denote the \textit{departments} and \textit{faculties} of educational institutions, for example.
\end{glyph}
\gindex{Course}{科}\strindex{9}{科}

\begin{comps}
\comprtwocone&禾 + 斗\\\hline
\comprthreecone&Grain (ear) + Dipper (Volume Measure)\\\hline
\end{comps}

\begin{remglyph}
\hooglig{The course} of \remcom{grain ears} is governed by the \remcom{Big Dipper}.\\\hline
\end{remglyph}

\begin{prons}
\pronsrtwocone&\textbf{カ (KA)}&---\\\hline
\pronsrthreecone&---&---\\\hline
\end{prons}

\begin{remprons}
\hooglig{A course} on how to stop \comon{cussing}.\\\hline
\end{remprons}

%{\middel}
% &  \mbox{()}\\\hline
\begin{somme}
科学&course + study = \textit{science}&カ{\middel}ガク \mbox{(KA{\middel}GAKU)}\\\hline
科目&course + eye = \textit{academic subject; course}&カ{\middel}モク \mbox{(KA{\middel}MOKU)}\\\hline
教科書&teach + course + write = \textit{textbook}&キョウ{\middel}カ{\middel}ショ \mbox{(KY\={O}{\middel}KA{\middel}SHO)}\\\hline
外科医&outside + course + doctor = \textit{surgeon}&ゲ{\middel}カ{\middel}イ \mbox{(GE{\middel}KA{\middel}I)}\\\hline
人文科学&humanity + science = \textit{humanities}&ジン{\middel}ブン{\middel}カ{\middel}ガク \mbox{(JIN{\middel}BUN{\middel}KA{\middel}GAKU)}\\\hline
医科&doctor + course = \textit{medical science/faculty}&イ{\middel}カ \mbox{(I{\middel}KA)}\\\hline
薬科&medicine + course = \textit{pharmacology}&ヤッ{\middel}カ \mbox{(YAK{\middel}KA)}\\\hline
前科&before + course = \textit{previous conviction}&ゼン{\middel}カ \mbox{(ZEN{\middel}KA)}\\\hline
\end{somme}

\begin{decomp}{5}
数学&は&面白い&科目&です。\\\hline
スウガク&は&おもしろい&カモク&です\\\hline
maths&は&interesting&subject&be/is\\\hline
\multicolumn{5}{|r|}{\textit{Mathematics is an interesting subject.}}\\\hline
\end{decomp}

\end{multicols}
%%--------------------------------------------------------------------
%4--------------------------------------------------------------------
\begin{multicols}{2}
\begin{glyph}{Younger Brother (弟)}{younger_brother}\label{glyph:younger_brother}
Younger brother.\\\hline
This 漢字 can also refer to the \textit{disciples} or \textit{pupils} of an accomplished master.\\\hline
Ensure that you can differentiate between this character and 第 (``No.'', from Entry \ref{glyph:number}, p. \pageref{glyph:number}).
\end{glyph}
\gindex{Younger Brother}{弟}\strindex{7}{弟}

\begin{comps}
\comprtwocone&艹 + 弔 + ノ\\\hline
\comprthreecone&Grass/Herb/Plant + Mourn (\ref{glyph:mourn}) + Slash/NO\\\hline
\end{comps}

\begin{remglyph}
\hooglig{My younger brother} \remcom{mourned} the \remcom{plants} that were \remcom{slashed}.\\\hline
\end{remglyph}

\begin{prons}
\pronsrtwocone&---&\textbf{おとうと (ot\={o}to)}\\\hline
\pronsrthreecone&テイ、ダイ、デ (TEI, DAI, DE)&---\\\hline
\end{prons}

\begin{remprons}
\hooglig{My younger brother} \ucomon{decorated} \comkun{a taut object} with \ucomon{dyed} \ucomon{tape}.\\\hline
\end{remprons}

%{\middel}
% &  \mbox{()}\\\hline
\begin{somme}
弟&\textit{younger brother}&おとうと \mbox{(ot\={o}to)}\\\hline
兄弟&elder brother + younger brother = \textit{brothers (and sisters)}&キョウ{\middel}ダイ \mbox{(KY\={O}{\middel}DAI)}\\\hline
弟子&younger brother + child = \textit{pupil; disciple}&デ{\middel}シ \mbox{(DE{\middel}SHI)}\\\hline
末弟&end + younger brother = \textit{youngest brother; last disciple}&マッ{\middel}テイ \mbox{(MAT{\middel}TEI)}\\\hline
兄弟分&brothers + part = \textit{buddy; sworn brother}&キョウ{\middel}ダイ{\middel}ブン \mbox{(KY\={O}{\middel}DAI{\middel}BUN)}\\\hline
兄弟愛&brothers + love = \textit{fraternal/brotherly love}&キョウ{\middel}ダイ{\middel}アイ \mbox{(KY\={O}{\middel}DAI{\middel}AI)}\\\hline
四海兄弟&whole world + brothers = \textit{universal brotherhood}&シ{\middel}カイ{\middel}ケイ{\middel}テイ \mbox{(SHI{\middel}KAI{\middel}KEI{\middel}TEI)}\\\hline
\end{somme}

\begin{decomp}{6}
彼&と&私&は&兄弟&です。\\\hline
かれ&と&わたし&は&キョウダイ&です\\\hline
he&と&I&は&brothers&be/is\\\hline
\multicolumn{6}{|r|}{\textit{He and I are brothers.}}\\\hline
\end{decomp}

\end{multicols}
\vspace*{\fill}
\pagebreak
%%--------------------------------------------------------------------
%5--------------------------------------------------------------------
\begin{multicols}{2}
\begin{glyph}{Manner (様)}{manner}\label{glyph:manner}
Manner.  Way.\\\hline
When following a family name, this 漢字 acts as a very polite form of \textit{Mr.}, \textit{Mrs.}, or \textit{Miss}.
\end{glyph}
\gindex{Manner}{様}\strindex{14}{様}

\begin{comps}
\comprtwocone & 木 + 羊 + 水 \\\hline
\comprthreecone & Tree + Sheep + Water \\\hline
\end{comps}

\begin{remglyph}
\hooglig{Manners} of \remcom{polite trees} give shade to \remcom{sheep} in search of \remcom{water}.\\\hline
\end{remglyph}

\begin{prons}
\pronsrtwocone&\textbf{ヨウ (Y\={O})}&\textbf{さま (sama)}\\\hline
\pronsrthreecone&---&---\\\hline
\end{prons}

\begin{remprons}
\hooglig{Mannerisms} of a \comon{yawning} \comkun{Samaritan}.\\\hline
\end{remprons}

%{\middel}
% &  \mbox{()}\\\hline
\begin{somme}
皆様&all + manner = \textit{Ladies and Gentlemen!}&みな{\middel}さま \mbox{(mina{\middel}sama)}\\\hline
神様&God + manner = \textit{God; god}&かみ{\middel}さま \mbox{(kami{\middel}sama)}\\\hline
様式&manner + type = \textit{mode; form}&ヨウ{\middel}シキ \mbox{(Y\={O}{\middel}SHIKI)}\\\hline
多様化&diverse + change = \textit{diversification}&タ{\middel}ヨウ{\middel}カ \mbox{(TA{\middel}Y\={O}{\middel}KA)}\\\hline
有様&have + manner = \textit{state; sight}&あ{\middel}り{\middel}さま \mbox{(a{\middel}ri{\middel}sama)}\\\hline
どの様&\textit{what sort?; what kind?}&どの{\middel}ヨウ \mbox{(dono{\middel}Y\={O})}\\\hline
左様なら&left + manner = \textit{good-bye}&サ{\middel}ヨウ{\middel}なら \mbox{(SA{\middel}Y\={O}{\middel}nara)}\\\hline
見様&see + manner = \textit{point of view}&み{\middel}ヨウ \mbox{mi{\middel}Y\={O}}\\\hline
様&\textit{plight; sad sight; mess; sorry state}&ざま \mbox{(zama)}\\\hline
\end{somme}

\begin{decomp}{6}
あなた&は&神様&を&信じます&か。\\\hline
あなた&は&かみさま&を&シンじます&か\\\hline
you&は&God&を&to believe&か\\\hline
\multicolumn{6}{|r|}{\textit{Do you believe in God?}}\\\hline
\end{decomp}

\end{multicols}
%%--------------------------------------------------------------------
%6--------------------------------------------------------------------
\begin{multicols}{2}
\begin{glyph}{Insect (虫)}{insect}\label{glyph:insect}
Insect.\\\hline
A secondary meaning is \textit{temper}.
\end{glyph}
\gindex{Insect}{虫}\strindex{6}{虫}

\begin{comps}
\comprtwocone&虫\\\hline
\comprthreecone&Insect\\\hline
\end{comps}

\begin{remglyph}
Part of the 漢字 radicals (\#{142}).\\\hline
\end{remglyph}

\begin{prons}
\pronsrtwocone&\textbf{チュウ (CH\={U})}&\textbf{むし (mushi)}\\\hline
\pronsrthreecone&---&---\\\hline
\end{prons}

\begin{remprons}
\hooglig{The insect} I \comon{chewed} on became all \comkun{mushi}.\\\hline
\end{remprons}

%{\middel}
% &  \mbox{()}\\\hline
\begin{somme}
虫&\textit{insect; bug}&むし \mbox{(mushi)}\\\hline
泣き虫&cry + insect = \textit{crybaby}&な{\middel}き{\middel}むし \mbox{(na{\middel}ki{\middel}mushi)}\\\hline
弱虫&weak + insect = \textit{coward; wimp}&よわ{\middel}むし \mbox{(yowa{\middel}mushi)}\\\hline
虫気&insect + spirit = \textit{nervous weakness}&むし{\middel}ケ \mbox{(mushi{\middel}KE)}\\\hline
玉虫&jewel + insect = \textit{jewel beetle}&たま{\middel}むし \mbox{(tama{\middel}mushi)}\\\hline
本の虫&book + insect = \textit{bookworm}&ホン{\middel}の{\middel}むし \mbox{(HON{\middel}no{\middel}mushi)}\\\hline
虫を殺す{\notcov}&temper + kill = \textit{to control one's temper}&むし{\middel}を{\middel}ころ{\middel}す \mbox{(mushi{\middel}wo{\middel}koro{\middel}su)}\\\hline
精虫{\notcov}&semen + insect = \textit{sperm}&セイ{\middel}チュウ \mbox{(SEI{\middel}CH\={U})}\\\hline
\notyet{成虫}&become + insect = \textit{adult insect}&セイ{\middel}チュウ \mbox{(SEI{\middel}CH\={U})}\\\hline
\end{somme}

\begin{decomp}{2}
虫も殺さぬ&顔。\\\hline
むしもころさぬ&かお\\\hline
innocent-looking&face\\\hline
\multicolumn{2}{|r|}{\textit{He looks as if he could not even harm a fly.}}\\\hline
\end{decomp}

\end{multicols}
\pagebreak
%%--------------------------------------------------------------------
%7--------------------------------------------------------------------
\begin{multicols}{2}
\begin{glyph}{Station (駅)}{station}\label{glyph:station}
Train station.
\end{glyph}
\gindex{Station}{駅}\strindex{14}{駅}

\begin{comps}
\comprtwocone&馬 + 尺\\\hline
\comprthreecone&Horse + Measurement (p. \pageref{glyph:measurement})\\\hline
\end{comps}

\begin{remglyph}
\hooglig{Stations} restrict \remcom{horses} above certain \remcom{measurements}.\\\hline
\end{remglyph}

\begin{prons}
\pronsrtwocone&\textbf{エキ (EKI)}&---\\\hline
\pronsrthreecone&---&---\\\hline
\end{prons}

\begin{remprons}
\hooglig{Station} an \comon{echinacea} greenhouse.\\\hline
{\warn}\textit{Echinacea} extract is the main ingredient in EchinaForce{\textregistered}, a herbal product for flu and colds.\\\hline
\end{remprons}

%{\middel}
% &  \mbox{()}\\\hline
\begin{somme}
駅&\textit{train station}&エキ \mbox{(EKI)}\\\hline
駅前&station + before = \textit{in front of the station}&エキ{\middel}まえ \mbox{(EKI{\middel}mae)}\\\hline
駅伝&station + transmit = \textit{long-distance relay race}&エキ{\middel}デン \mbox{(EKI{\middel}DEN)}\\\hline
各駅&every + station = \textit{every station}&カク{\middel}エキ \mbox{(KAKU{\middel}EKI)}\\\hline
主要駅&main + station = \textit{principal stations}&シュ{\middel}ヨウ{\middel}エキ \mbox{(SHU{\middel}Y\={O}{\middel}EKI)}\\\hline
着駅&wear + station = \textit{arriving station}&チャク{\middel}エキ \mbox{(CHAKU{\middel}EKI)}\\\hline
地下駅&below ground + station = \textit{underground station}&チ{\middel}カ{\middel}エキ \mbox{(CHI{\middel}KA{\middel}EKI)}\\\hline
自動車駅&automobile + station = \textit{bus terminal}&ジ{\middel}ドウ{\middel}シャ{\middel}エキ \mbox{(JI{\middel}D\={O}{\middel}SHA{\middel}EKI)}\\\hline
\end{somme}

\begin{decomp}{5}
私&は&駅&へ&行った。\\\hline
わたし&は&エキ&え&いった\\\hline
I&は&station&え&to go\\\hline
\multicolumn{5}{|r|}{\textit{I went to the station.}}\\\hline
\end{decomp}

\end{multicols}
%%--------------------------------------------------------------------
%8--------------------------------------------------------------------
\begin{multicols}{2}
\begin{glyph}{Commerce (商)}{commerce}\label{glyph:commerce}
Commerce.  Trade.
\end{glyph}
\gindex{Commerce}{商}\strindex{11}{商}

\begin{comps}
\comprtwocone&立 + 冏\\\hline
\comprthreecone&Stand + Bright\\\hline
\rye{2}{冏 = 冂 + 儿 + 口.  Bright = Wilderness + (Person) Legs + Mouth.  \hooglig{Bright} \remcom{people} \remcom{enclose} \remcom{vampires} in the \remcom{wilderness}.}
\end{comps}

\begin{remglyph}
\hooglig{Commerce} took a \remcom{bright} \remcom{stand}.\\\hline
\end{remglyph}

\begin{prons}
\pronsrtwocone&\textbf{ショウ (SH\={O})}&---\\\hline
\pronsrthreecone&---&あきな (akina)\\\hline
\end{prons}

\begin{remprons}
\hooglig{Commerce} is \ucomkun{akin} to a \comon{shore}; it comes and goes.\\\hline
\end{remprons}

%{\middel}
% &  \mbox{()}\\\hline
\begin{somme}
商人&commerce + person = \textit{merchant; dealer}&ショウ{\middel}ニン \mbox{(SH\={O}{\middel}NIN)}\\\hline
商売&commerce + sell = \textit{commerce; trade}&ショウ{\middel}バイ \mbox{(SH\={O}{\middel}BAI)}\\\hline
通商&passage + commerce = \textit{commerce; trade}&ツウ{\middel}ショウ \mbox{(TS\={U}{\middel}SH\={O})}\\\hline
商店&commerce + shop = \textit{shop; store}&ショウ{\middel}テン \mbox{(SH\={O}{\middel}TEN)}\\\hline
商品&commerce + goods = \textit{articles of commerce}&ショウ{\middel}ヒン \mbox{(SH\={O}{\middel}HIN)}\\\hline
年商&year + commerce = \textit{yearly turnover}&ネン{\middel}ショウ \mbox{(NEN{\middel}SH\={O})}\\\hline
月商&moon + commerce = \textit{monthly sales}&ゲッ{\middel}ショウ \mbox{(GES{\middel}SH\={O})}\\\hline
水商売&water + commerce = \textit{chancy business}&みず{\middel}ショウ{\middel}バイ \mbox{(mizu{\middel}SH\={O}{\middel}BAI)}\\\hline
商取引&commerce + transactions = \textit{business transaction}&ショウ{\middel}と{\middel}り{\middel}ひ{\middel}き \mbox{(SH\={O}{\middel}to{\middel}ri{\middel}hi{\middel}ki)}\\\hline
\end{somme}

\begin{decomp}{5}
近く&に&商店街&が&ある。\\\hline
ちかく&に&ショウテンガイ&が&ある\\\hline
near&に&shopping district&が&to be\\\hline
\multicolumn{5}{|r|}{\textit{There is a shopping area nearby.}}\\\hline
\end{decomp}

\end{multicols}
\pagebreak
%%--------------------------------------------------------------------
%9--------------------------------------------------------------------
\begin{multicols}{2}
\begin{glyph}{Bad (悪)}{bad}\label{glyph:bad}
Bad.  Wrong.  Inferior\ldots\\\hline
\end{glyph}
\gindex{Bad}{悪}\strindex{11}{悪}

\begin{comps}
\comprtwocone&亜 + 心\\\hline
\comprthreecone&Asia (p. \pageref{glyph:asia}) + Heart\\\hline
\end{comps}

\begin{remglyph}
\hooglig{Bad} \remcom{asian} \remcom{hearts}.\\\hline
\end{remglyph}

\begin{prons}
\pronsrtwocone&\textbf{アク (AKU)}&\textbf{わる (waru)}\\\hline
\pronsrthreecone&オ (O)&---\\\hline
\end{prons}

\begin{remprons}
\hooglig{Bad} \comkun{war} tactics made \comon{Aku} \ucomon{object}.\\\hline
{\warn}\textit{Aku} is a fictional character from \textit{Samurai Jack}.\\\hline
\end{remprons}

%{\middel}
% &  \mbox{()}\\\hline
\begin{somme}
悪い&\textit{bad}&わる{\middel}い \mbox{(waru{\middel}i)}\\\hline
最悪&most + bad = \textit{worst}&サイ{\middel}アク \mbox{(SAI{\middel}AKU)}\\\hline
悪意&bad + mind = \textit{evil intention}&アク{\middel}イ \mbox{(AKU{\middel}I)}\\\hline
悪事&bad + matter = \textit{crime; vice}&アク{\middel}ジ \mbox{(AKU{\middel}JI)}\\\hline
{\diff}意地悪&disposition + bad = \textit{malicious; unkind}&イ{\middel}ジ{\middel}わる \mbox{(I{\middel}JI{\middel}waru)}\\\hline
悪化&bad + change = \textit{degeneration; corruption}&アッ{\middel}カ \mbox{(AK{\middel}KA)}\\\hline
凶悪&evil + bad = \textit{atrocious; brutal; villainous}&キョウ{\middel}アク \mbox{(KY\={O}{\middel}AKU)}\\\hline
気持ち悪い&feeling + bad = \textit{bad feeling; unpleasant; revolting}&キ{\middel}も{\middel}ち{\middel}わる{\middel}い \mbox{(KI{\middel}mo{\middel}chi{\middel}waru{\middel}i)}\\\hline
\end{somme}

\begin{decomp}{3}
気分&悪い&の？\\\hline
キブン&わるい&の\\\hline
feeling/mood&bad&の\\\hline
\multicolumn{3}{|r|}{\textit{Are you feeling sick?}}\\\hline
\end{decomp}

\end{multicols}
%%--------------------------------------------------------------------
%10-------------------------------------------------------------------
\begin{multicols}{2}
\begin{glyph}{Know (知)}{know}\label{glyph:know}
Know.  Information.
\end{glyph}
\gindex{Know}{知}\strindex{8}{知}

\begin{comps}
\comprtwocone&矢 + 口\\\hline
\comprthreecone&Arrow + Mouth\\\hline
\end{comps}

\begin{remglyph}
\hooglig{Knowing information} will equip you with \remcom{arrows} in your \remcom{mouth}, allowing you to make witty comments.\\\hline
\end{remglyph}

\begin{prons}
\pronsrtwocone & \textbf{チ (CHI)} & \textbf{し (shi)}\\\hline
\pronsrthreecone & --- & --- \\\hline
\end{prons}

\begin{remprons}
\hooglig{Know} \comon{cheap} \comkun{shifting} techniques.\\\hline
\end{remprons}

%{\middel}
% &  \mbox{()}\\\hline
\begin{somme}
知る&\textit{know}&し{\middel}る \mbox{(shi{\middel}ru)}\\\hline
思い知る&think + know = \textit{to realise}&おも{\middel}い{\middel}し{\middel}る \mbox{(omo{\middel}i{\middel}shi{\middel}ru)}\\\hline
知らせ&\textit{notice; information}&し{\middel}らせ \mbox{(shi{\middel}rase)}\\\hline
知人&know + person = \textit{acquaintance}&チ{\middel}ジン \mbox{(CHI{\middel}JIN)}\\\hline
通知&passage + know = \textit{notice; notification; report; posting}&ツウ{\middel}チ \mbox{(TS\={U}{\middel}CHI)}\\\hline
未知&not yet + know = \textit{not yet known; unknown}&ミ{\middel}チ \mbox{(MI{\middel}CHI)}\\\hline
予知&beforehand + know = \textit{foreknowledge}&ヨ{\middel}チ \mbox{(YO{\middel}CHI)}\\\hline
理知的&intellect + target = \textit{intellectual}&リ{\middel}チ{\middel}テキ \mbox{(RI{\middel}CHI{\middel}TEKI)}\\\hline
良知&good + know = \textit{intuition}&リョウ{\middel}チ \mbox{(RY\={O}{\middel}CHI)}\\\hline
\end{somme}
\end{multicols}

\begin{decomp}{10}
彼&に&は&知人&は&多い&が&友人&は&少ない。\\\hline
かれ&に&は&チジン&は&おおい&が&ユウジン&は&すくない\\\hline
he&に&は&acquaintances&は&many&が&friend&は&few\\\hline
\multicolumn{10}{|r|}{\textit{He has many acquaintances but few friends.}}\\\hline
\end{decomp}
\pagebreak
%%--------------------------------------------------------------------
%11-------------------------------------------------------------------
\begin{multicols}{2}
\begin{glyph}{Borrow (借)}{borrow}\label{glyph:borrow}
Borrow.
\end{glyph}
\gindex{Borrow}{借}\strindex{10}{借}

\begin{comps}
\comprtwocone&人亻 + 昔\\\hline
\comprthreecone&Person + Long Ago (p. \pageref{glyph:long_ago})\\\hline
\end{comps}

\begin{remglyph}
\hooglig{Borrowed} from a \remcom{person} \remcom{long ago}.\\\hline
\end{remglyph}

\begin{prons}
\pronsrtwocone&\textbf{シャク (SHAKU)}&\textbf{か (ka)}\\\hline
\pronsrthreecone&---&---\\\hline
\end{prons}

\begin{remprons}
\hooglig{Borrowing} \comon{shucks} is \comkun{customary}.\\\hline
{\warn}\textit{Shuck} refers to an outer covering such as a husk or pod.\\\hline
\end{remprons}

%{\middel}
% &  \mbox{()}\\\hline
\begin{somme}
\rye{3}{The verb form of this 漢字 can incorporate its ひらがな ending when forming nouns.}
借りる&\textit{borrow}&か{\middel}りる \mbox{(ka{\middel}riru)}\\\hline
借金&borrow + money = \textit{loan; a debt}&シャッ{\middel}キン \mbox{(SHAK{\middel}KIN)}\\\hline
借入金&debt + money = \textit{debt}&かり{\middel}いれ{\middel}キン \mbox{(kari{\middel}ire{\middel}KIN)}\\\hline
借家&borrow + house = \textit{house for rent}&シャク{\middel}や \mbox{(SHAKU{\middel}ya)}\\\hline
間借り&interval + borrow = \textit{renting a room}&ま{\middel}が{\middel}り \mbox{(ma{\middel}ga{\middel}ri)}\\\hline
借り手&borrow + hand = \textit{borrower; debtor; tenant}&か{\middel}り{\middel}て \mbox{(ka{\middel}ri{\middel}te)}\\\hline
貸し借り&lend + borrow = \textit{lending and borrowing}&か{\middel}し{\middel}か{\middel}り \mbox{(ka{\middel}shi{\middel}ka{\middel}ri)}\\\hline
\end{somme}
\end{multicols}

\begin{decomp}{6}
私&に&は&借金&が&無い。\\\hline
わたし&に&は&シャッキン&が&ない\\\hline
I&に&は&debt&が&non-existent\\\hline
\multicolumn{6}{|r|}{\textit{I am free of debt.}}\\\hline
\end{decomp}
%%--------------------------------------------------------------------
%12-------------------------------------------------------------------
\begin{multicols}{2}
\begin{glyph}{Decide (決)}{decide}\label{glyph:decide}
Decide.
\end{glyph}
\gindex{Decide}{決}\strindex{7}{決}

\begin{comps}
\comprtwocone&水氵氺 + ユ + 人\\\hline
\comprthreecone&Water + Hangman/Gallows/YU + person\\\hline
\end{comps}

\begin{remglyph}
\hooglig{Decide} between drowning in \remcom{water}, or being \remcom{hanged} in front of \remcom{people}.\\\hline
\end{remglyph}

\begin{prons}
\pronsrtwocone & \textbf{ケツ (KETSU)} & \textbf{き (ki)}\\\hline
\pronsrthreecone & --- & --- \\\hline
\end{prons}

\begin{remprons}
\hooglig{Decide} to \comkun{keep} the \comon{Kets} alive.\\\hline
\end{remprons}

%{\middel}
% &  \mbox{()}\\\hline
\begin{somme}
決まる (intr.)&\textit{to be decided / settled}&き{\middel}まる \mbox{(ki{\middel}maru)}\\\hline
決める (tr.)&\textit{to decide / settle}&き{\middel}める \mbox{(ki{\middel}meru)}\\\hline
決心&decide + heart = \textit{determination; resolution}&ケッ{\middel}シン \mbox{(KES{\middel}SHIN)}\\\hline
決行&decide + go = \textit{doing with resolve}&ケッ{\middel}コウ \mbox{(KEK{\middel}K\={O})}\\\hline
決着&decide + wear = \textit{conclusion; end}&ケッ{\middel}チャク \mbox{(KEC{\middel}CHAKU)}\\\hline
意思決定&intention + decision = \textit{decision-making}&イ{\middel}シ{\middel}ケッ{\middel}テイ \mbox{(I{\middel}SHI{\middel}KET{\middel}TEI)}\\\hline
決別&decide + separate = \textit{farewell; parting}&ケツ{\middel}ベツ \mbox{(KETSU{\middel}BETSU)}\\\hline
\end{somme}

\end{multicols}

\begin{decomp}{6}
今&こそ&決心す&べき&時&だ。\\\hline
いま&こそ&ケッシンす&べき&とき&だ\\\hline
now&for sure&determination&shall/should/must&time&be/is\\\hline
\multicolumn{6}{|r|}{\textit{Now is when you have to make up your mind.}}\\\hline
\end{decomp}
\pagebreak
%%--------------------------------------------------------------------
%13-------------------------------------------------------------------
\begin{multicols}{2}
\begin{glyph}{Partition (区)}{partition}\label{glyph:partition}
Partition.  Division.\\\hline
This character is used to indicate \textit{zones} or \textit{areas}.  You will often see this 漢字 in mailing addresses, as it's used to denote a \textit{ward} or \textit{district} of a city.
\end{glyph}
\gindex{Partition}{区}\strindex{4}{区}

\begin{comps}
\comprtwocone&匚 + 乂\\\hline
\comprthreecone&Square Box + NO/Slash + Dot\\\hline
\end{comps}

\begin{remglyph}
\hooglig{Partition} \remcom{square box} with \remcom{crossed swords}.\\\hline
\end{remglyph}

\begin{prons}
\pronsrtwocone&\textbf{ク (KU)}&---\\\hline
\pronsrthreecone&---&---\\\hline
\end{prons}

\begin{remprons}
\hooglig{The division} for \comon{cooking}.\\\hline
\end{remprons}

%{\middel}
% &  \mbox{()}\\\hline
\begin{somme}
区分&partition + part = \textit{partition; division}&ク{\middel}ブン \mbox{(KU{\middel}BUN)}\\\hline
区別&partition + different = \textit{distinction; difference}&ク{\middel}ベツ \mbox{(KU{\middel}BETSU)}\\\hline
新宿区&new + lodging + partition = \textit{Shinjuku ward (in Tokyo)}&シン{\middel}ジュク{\middel}ク \mbox{(SHIN{\middel}JUKU{\middel}KU)}\\\hline
学区&study + partition = \textit{school district}&ガッ{\middel}ク \mbox{(GAK{\middel}KU)}\\\hline
工区&craft + partition = \textit{construction area}&コウ{\middel}ク \mbox{(K\={O}{\middel}KU)}\\\hline
線区&line + partition = \textit{train line; bus line}&セン{\middel}ク \mbox{(SEN{\middel}KU)}\\\hline
教区&teach + partition = \textit{parish}&キョウ{\middel}ク \mbox{(KY\={O}{\middel}KU)}\\\hline
\end{somme}
\end{multicols}

\begin{decomp}{10}
犬&は&白&と&黒&と&の&区別&が&付く。\\\hline
いぬ&は&しろ&と&くろ&と&の&クベツ&が&つく\\\hline
dog&あ&white&と&black&と&の&distinction&が&to be attached\\\hline
\multicolumn{10}{|r|}{\textit{The dog knows black from white.}}\\\hline
\end{decomp}
%%--------------------------------------------------------------------
%14-------------------------------------------------------------------
\begin{multicols}{2}
\begin{glyph}{Discharge (発)}{discharge}\label{glyph:discharge}
Discharge.  Break out.
\end{glyph}
\gindex{Discharge}{発}\strindex{9}{発}

\begin{comps}
\comprtwocone&癶 +　ニ　+ 儿\\\hline
\comprthreecone&Footsteps/Departure + Two + Man/Person\\\hline
\end{comps}

\begin{remglyph}
\hooglig{Discharge} \remcom{both} \remcom{people} for their \remcom{departure}.\\\hline
\hooglig{Discharge} the \remcom{departing footsteps} of \remcom{two} \remcom{human legs}.\\\hline
\end{remglyph}

\begin{prons}
\pronsrtwocone&\textbf{ハツ (HATSU)}&---\\\hline
\pronsrthreecone&ホツ (HOTSU)&---\\\hline
\end{prons}

\begin{remprons}
\hooglig{Break} out of the \comon{huts} to reach the \ucomon{hotspot}.\\\hline
\end{remprons}

%{\middel}
% &  \mbox{()}\\\hline
\begin{somme}
発音&discharge + sound = \textit{pronunciation}&ハツ{\middel}オン \mbox{(HATSU{\middel}ON)}\\\hline
出発&exit + discharge = \textit{departure}&シュッ{\middel}パツ \mbox{(SHUP{\middel}PATSU)}\\\hline
新発売&new + discharge + sell = \textit{``Now on Sale''}&シン{\middel}ハツ{\middel}バイ \mbox{(SHIN{\middel}HATSU{\middel}BAI)}\\\hline
発明&discharge + bright = \textit{invention}&ハツ{\middel}メイ \mbox{(HATSU{\middel}MEI)}\\\hline
発表&discharge + express = \textit{announcement{\ldots}}&ハッ{\middel}ピョウ \mbox{(HAP{\middel}PY\={O})}\\\hline
発見&discharge + see = \textit{discovery; finding}&ハッ{\middel}ケン \mbox{(HAK{\middel}KEN)}\\\hline
発足&discharge + leg = \textit{starting; launch; founding}&ハッ{\middel}ソク \mbox{(HAS{\middel}SOKU)}\\\hline
発作&discharge + make = \textit{fit; spasm; seizure}&ハッ{\middel}サ \mbox{(HAS{\middel}SA)}\\\hline
多発&many + discharge = \textit{repeated occurrence}&タ{\middel}ハツ \mbox{(TA{\middel}HATSU)}\\\hline
\end{somme}

\begin{decomp}{3}
出発する&時間&だ。\\\hline
シュッパツする&ジカン&だ\\\hline
departure&time&be/is\\\hline
\multicolumn{3}{|r|}{\textit{It is really time for us to go.}}\\\hline
\end{decomp}

\end{multicols}
\pagebreak
%%--------------------------------------------------------------------
%15-------------------------------------------------------------------
\begin{multicols}{2}
\begin{glyph}{How Many (幾)}{how_many}\label{glyph:how_many}
How many.  How much.\\\hline
Another meaning is conveyed when this character precedes a `number' 漢字; in these cases the character means \textit{some\ldots} or \textit{several\ldots}
\end{glyph}
\gindex{How Many}{幾}\strindex{12}{幾}

\begin{comps}
\comprtwocone&戍 + 幺\\\hline
\comprthreecone&Protection + Short Thread\\\hline
\rye{2}{戍 = 戈 + 人.  Protection = Halberd + Person.  \hooglig{Protect} the \remcom{person} with a \remcom{halberd}.}
\end{comps}

\begin{remglyph}
\hooglig{How many} \remcom{shorties} do you \remcom{protect}?\\\hline
\end{remglyph}

\begin{prons}
\pronsrtwocone&---&\textbf{いく (iku)}\\\hline
\pronsrthreecone&キ (KI)&---\\\hline
\end{prons}

\begin{remprons}
\hooglig{How many} \ucomon{killings} are you going to \comkun{eke out}?\\\hline
\end{remprons}

%{\middel}
% &  \mbox{()}\\\hline
\begin{somme}
幾人&how many + person = \textit{number of people}&いく{\middel}ニン \mbox{(iku{\middel}NIN)}\\\hline
幾つ&\textit{how much; how many}&いく{\middel}つ \mbox{(iku{\middel}tsu)}\\\hline
幾ら&\textit{how much}&いく{\middel}ら \mbox{(iku{\middel}ra)}\\\hline
幾百&how much + hundred = \textit{some/several hundreds}&いく{\middel}ヒャク \mbox{(iku{\middel}HYAKU)}\\\hline
幾日&how much + day = \textit{how many days; what day of the month}&いく{\middel}ニチ \mbox{(iku{\middel}NICHI)}\\\hline
幾多&how many + many = \textit{many; numerous}&いくタ \mbox{(iku{\middel}TA)}\\\hline
幾重&how many + heavy = \textit{piling up; many layers}&いく{\middel}え \mbox{(iku{\middel}e)}\\\hline
\notyet{平面幾何学}&level surface + geometry = \textit{plane geometry}&ヘイ{\middel}メン{\middel}キ{\middel}カ{\middel}ガク (HEI{\middel}MEN{\middel}KI{\middel}KA{\middel}GAKU)\\\hline
\end{somme}

\begin{decomp}{4}
幾ら&で&買う&か。\\\hline
いくら&で&かう&か\\\hline
how many&で&to buy&か\\\hline
\multicolumn{4}{|r|}{\textit{What price will you give for this?}}\\\hline
\end{decomp}

\end{multicols}
%%--------------------------------------------------------------------
%16-------------------------------------------------------------------
\begin{multicols}{2}
\begin{glyph}{Room (室)}{room}\label{glyph:room}
Room.  Chamber.\\\hline
Don't confuse this character with 屋, the 漢字 for \textit{Roof} (Entry \ref{glyph:roof}, p. \pageref{glyph:roof}).
\end{glyph}
\gindex{Room}{室}\strindex{9}{室}

\begin{comps}
\comprtwocone&宀 + 至\\\hline
\comprthreecone&Roof + Reach (p. \pageref{glyph:reach})\\\hline
\end{comps}

\begin{remglyph}
\hooglig{The room's} \remcom{roof} is \remcom{unreachable}.\\\hline
\end{remglyph}

\begin{prons}
\pronsrtwocone&\textbf{シツ (SHITSU)}&---\\\hline
\pronsrthreecone&---&むろ (muro)\\\hline
\end{prons}

\begin{remprons}
\hooglig{The room} had \ucomkun{muron} under the \comon{sheets}.\\\hline
{\warn}\textit{Muron} is a rice cake in Filipino cuisine.\\\hline
\end{remprons}

%{\middel}
% &  \mbox{()}\\\hline
\begin{somme}
室内楽&indoors + music = \textit{chamber music}&シツ{\middel}ナイ{\middel}ガク \mbox{(SHITSU{\middel}NAI{\middel}GAKU)}\\\hline
室外&room + outside = \textit{outdoor(s)}&シツ{\middel}ガイ \mbox{(SHITSU{\middel}GAI)}\\\hline
待合室&wait + match + room = \textit{waiting room}&まち{\middel}あい{\middel}シツ \mbox{(machi{\middel}ai{\middel}SHITSU)}\\\hline
地下室&underground + room = \textit{cellar; basement}&チ{\middel}カ{\middel}シツ \mbox{(CHI{\middel}KA{\middel}SHITSU)}\\\hline
和室&harmony + rom = \textit{Japanese-style room}&ワ{\middel}シツ \mbox{(WA{\middel}SHITSU)}\\\hline
読書室&reading + room = \textit{reading room}&ドク{\middel}ショ{\middel}シツ (DOKU{\middel}SHO{\middel}SHITSU)\\\hline
工作室&handicraft + room = \textit{workshop}&コウ{\middel}サク{\middel}シツ \mbox{(K\={O}{\middel}SAKU{\middel}SHITSU)}\\\hline
\notyet{温室}&warm + room = \textit{hothouse; greenhouse}&オン{\middel}シツ \mbox{(ON{\middel}SHITSU)}\\\hline
\end{somme}

\begin{decomp}{4}
教室&で&食べる&の。\\\hline
キョウシツ&で&たべる&の\\\hline
classroom&で&to eat&の\\\hline
\multicolumn{4}{|r|}{\textit{I eat in the classroom.}}\\\hline
\end{decomp}

\end{multicols}
\pagebreak
%%--------------------------------------------------------------------
%17-------------------------------------------------------------------
\begin{multicols}{2}
\begin{glyph}{Fall (落)}{fall}\label{glyph:fall}
Fall.  Fail.
\end{glyph}
\gindex{Fall}{落}\strindex{12}{落}

\begin{comps}
\comprtwocone&艹 + 洛\\\hline
\comprthreecone&Grass/Herb/Plant + (Kyoto = Water + Each)\\\hline
\rye{2}{洛 = 水氵氺 + 各.  Ky\={o}to = Water + Each/Every.  \hooglig{Ky\={o}to} has \remcom{water} \remcom{each} time.}
\end{comps}

\begin{remglyph}
\hooglig{Falling} \remcom{grass blades} is a beautiful sight in \remcom{Ky\={o}to}.\\\hline
\end{remglyph}

\begin{prons}
\pronsrtwocone&\textbf{ラク (RAKU)}&\textbf{お (o)}\\\hline
\pronsrthreecone&---&---\\\hline
\end{prons}

\begin{remprons}
\hooglig{Falling} was \comon{luckily} not \comkun{obligatory}.\\\hline
\end{remprons}

%{\middel}
% &  \mbox{()}\\\hline
\begin{somme}
\rye{3}{The verbs below can sometimes absorb their ひらがな endings when forming words with other 漢字.}
落ちる (intr.)&\textit{to fall}&お{\middel}ちる \mbox{(o{\middel}chiru)}\\\hline
落とす (tr.)&\textit{to fall/drop (sth.)}&お{\middel}とす \mbox{(o{\middel}tosu)}\\\hline
急落&hasten + fall = \textit{sudden drop/decline}&キュウ{\middel}ラク \mbox{(KY\={U}{\middel}RAKU)}\\\hline
落ち着く&fall + wear = \textit{to calm down; to compose oneself}&お{\middel}ち{\middel}つ{\middel}く \mbox{(o{\middel}chi{\middel}tsu{\middel}ku)}\\\hline
落とし物&fall + thing = \textit{lost property}&お{\middel}とし{\middel}もの \mbox{(o{\middel}toshi{\middel}mono)}\\\hline
集落&gather + fall = \textit{settlement; community}&シュウ{\middel}ラク \mbox{(SH\={U}{\middel}RAKU)}\\\hline
落書き&fall + write = \textit{scribble; graffiti}&ラク{\middel}が{\middel}き \mbox{(RAKU{\middel}ga{\middel}ki)}\\\hline
落雷&fall + thunder = \textit{lightning strike; thunderbolt}&ラク{\middel}ライ \mbox{(RAKU{\middel}RAI)}\\\hline
\end{somme}

\begin{decomp}{3}
葉&が&落ちた。\\\hline
は&が&おちた\\\hline
leaf&が&to fall down\\\hline
\multicolumn{3}{|r|}{\textit{The leaves fell.}}\\\hline
\end{decomp}

\end{multicols}
%%--------------------------------------------------------------------
%18-------------------------------------------------------------------
\begin{multicols}{2}
\begin{glyph}{Fast (速)}{fast}\label{glyph:fast}
Fast, in terms of speed and movement.  Quick.\\\hline
Recall that 早 (\textit{Early}, Entry \ref{glyph:early}, p. \pageref{glyph:early}) is related to the idea of \textit{fast} in terms of time.
\end{glyph}
\gindex{Fast}{速}\strindex{10}{速}

\begin{comps}
\comprtwocone&束 + 辶\\\hline
\comprthreecone&Bundle (p. \pageref{glyph:bundle}) + Road/Walk\\\hline
\end{comps}

\begin{remglyph}
\hooglig{Fast} \remcom{bundles} of people on the \remcom{road}.\\\hline
\end{remglyph}

\begin{prons}
\pronsrtwocone&\textbf{ソク (SOKU)}&\textbf{はや (haya)}\\\hline
\pronsrthreecone&---&すみ (sumi)\\\hline
\end{prons}

\begin{remprons}
\hooglig{The fast} \comon{sock} and \comkun{hay allergy} is what they are \ucomkun{sueing me} for.\\\hline
\end{remprons}

%{\middel}
% &  \mbox{()}\\\hline
\begin{somme}
速い&\textit{fast}&はや{\middel}い \mbox{(haya{\middel}i)}\\\hline
速く&\textit{quickly; swiftly}&はや{\middel}く \mbox{(haya{\middel}ku)}\\\hline
高速&tall + fast = \textit{high-speed}&コウ{\middel}ソク \mbox{(K\={O}{\middel}SOKU)}\\\hline
速度&fast + degree = \textit{speed}&ソク{\middel}ド \mbox{(SOKU{\middel}DO)}\\\hline
早速&early + fast = \textit{at once; immediately}&サッ{\middel}ソク \mbox{(SAS{\middel}SOKU)}\\\hline
秒速&second + fast = \textit{per second}&ビョウ{\middel}ソク \mbox{(BY\={O}{\middel}SOKU)}\\\hline
全速力&complete + speed = \textit{full speed}&ゼン{\middel}ソク{\middel}リョク \mbox{(ZEN{\middel}SOKU{\middel}RYOKU)}\\\hline
角速度&angle + speed = \textit{angular velocity}&カク{\middel}ソク{\middel}ド \mbox{(KAKU{\middel}SOKU{\middel}DO)}\\\hline
最速化&fastest + change = \textit{optimisation}&サイ{\middel}ソク{\middel}カ \mbox{(SAI{\middel}SOKU{\middel}KA)}\\\hline
変速&change + speed = \textit{shifting gears}&ヘン{\middel}ソク \mbox{(HEN{\middel}SOKU)}\\\hline
\end{somme}

\begin{decomp}{3}
列車&は&加速した。\\\hline
レッシャ&は&カソクした\\\hline
train&は&speeding up\\\hline
\multicolumn{3}{|r|}{\textit{The train gathered speed.}}\\\hline
\end{decomp}

\end{multicols}
\pagebreak
%%--------------------------------------------------------------------
%19-------------------------------------------------------------------
\begin{multicols}{2}
\begin{glyph}{Color (色)}{color}\label{glyph:color}
Color.\\\hline
Interestingly, this 漢字 also conveys a decidedly erotic meaning to many compounds in which it appears.
\end{glyph}
\gindex{Color}{色}\strindex{6}{色}

\begin{comps}
\comprtwocone&色\\\hline
\comprthreecone&Color\\\hline
\end{comps}

\begin{remglyph}
Part of the 漢字 radicals (\#{139}).\\\hline
\end{remglyph}

\begin{prons}
\pronsrtwocone&\textbf{ショク (SHOKU)}&\textbf{いろ (iro)}\\\hline
\rye{3}{いろ is more common in the first position, and ショク in the second and third.}
\pronsrthreecone&シキ (SHIKI)&---\\\hline
\end{prons}

\begin{remprons}
\hooglig{The colour} of the \comkun{iroko} tree \comon{shocked} the \ucomon{Sheikh}.\\\hline
{\warn}\textit{Iroko} is a large hardwood tree from the west coast of tropical Africa that can live up to 500 years.\\\hline
\end{remprons}

%{\middel}
% &  \mbox{()}\\\hline
\begin{somme}
色目&color + eye = \textit{amorous glance}&いろ{\middel}め \mbox{(iro{\middel}me)}\\\hline
原色&original + color = \textit{primary color}&ゲン{\middel}ショク \mbox{(GEN{\middel}SHOKU)}\\\hline
特色&special + color = \textit{characteristic}&トク{\middel}ショク \mbox{(TOKU{\middel}SHOKU)}\\\hline
顔色&face + color = \textit{countenance}&かお{\middel}いろ \mbox{(kao{\middel}iro)}\\\hline
好色&like + color = \textit{lasciviousness}&コウ{\middel}ショク \mbox{K\={O}{\middel}SHOKU}\\\hline
物色&thing + color = \textit{looking for sth.}&ブッ{\middel}ショク \mbox{(BUS{\middel}SHOKU)}\\\hline
出色&exit + color = \textit{excellent}&シュッ{\middel}ョク \mbox{(SHUS{\middel}SHOKU)}\\\hline
色道&color + road = \textit{sexual passion}&シキ{\middel}ドウ \mbox{(SHIKI{\middel}D\={O})}\\\hline
色町&color + town = \textit{red-light district}&いろ{\middel}まち \mbox{(iro{\middel}machi)}\\\hline
\end{somme}

\begin{decomp}{2}
色っぽい&女性。\\\hline
いろっぽい&ジョセイ\\\hline
sexy&woman\\\hline
\multicolumn{2}{|r|}{\textit{She is a fox.}}\\\hline
\end{decomp}

\end{multicols}
%%--------------------------------------------------------------------
%20-------------------------------------------------------------------
\begin{multicols}{2}
\begin{glyph}{Start (始)}{start}\label{glyph:start}
Start.
\end{glyph}
\gindex{Start}{始}\strindex{8}{始}

\begin{comps}
\comprtwocone&女 + 台\\\hline
\comprthreecone& Woman + Platform (p. \pageref{glyph:platform})\\\hline
\end{comps}

\begin{remglyph}
\hooglig{The start} of any \remcom{woman} is some \remcom{platform} they use.\\\hline
\end{remglyph}

\begin{prons}
\pronsrtwocone&\textbf{シ (SHI)}&\textbf{はじ (haji)}\\\hline
\pronsrthreecone&---&---\\\hline
\end{prons}

\begin{remprons}
\hooglig{Start} the \comkun{hajj} to \comon{shield} you from evil forces.\\\hline
{\warn}\textit{The Hajj} is an annual Islamic pilgrimage to Mecca.\\\hline
\end{remprons}

%{\middel}
% &  \mbox{()}\\\hline
\begin{somme}
始まる (intr.)&\textit{to start}&はじ{\middel}まる \mbox{(haji{\middel}maru)}\\\hline
始める (tr.)&\textit{to start (sth.)}&はじ{\middel}める \mbox{(haji{\middel}meru)}\\\hline
原始&original + start = \textit{primitive}&ゲン{\middel}シ \mbox{(GEN{\middel}SHI)}\\\hline
始発&start + discharge = \textit{first departure}&シ{\middel}ハツ \mbox{(SHI{\middel}HATSU)}\\\hline
始動&start + move = \textit{activation}&シ{\middel}ドウ \mbox{(SHI{\middel}D\={O})}\\\hline
事始め&matter + start = \textit{start of new things}&こと{\middel}はじ{\middel}め \mbox{(koto{\middel}haji{\middel}me)}\\\hline
更始&anew + start = \textit{reform; renewal}&コウ{\middel}シ \mbox{(K\={O}{\middel}SHI)}\\\hline
開始日&open + start + day = \textit{start date}&カイ{\middel}シ{\middel}び \mbox{(KAI{\middel}SHI{\middel}bi)}\\\hline
\notyet{末始終}&end + from beginning to end = \textit{forever; for life}&すえ{\middel}シ{\middel}ジュウ \mbox{(sue{\middel}SHI{\middel}J\={U})}\\\hline
\end{somme}

\begin{decomp}{5}
何&時&に&始まる&の？\\\hline
なん&ジ&に&はじまる&の\\\hline
what&time&に&to start&の\\\hline
\multicolumn{5}{|r|}{\textit{What time does it start?}}\\\hline
\end{decomp}

\end{multicols}
\pagebreak
%%--------------------------------------------------------------------
%21-------------------------------------------------------------------
\vspace*{\fill}
\begin{multicols}{2}
\begin{glyph}{Approve (是)}{approve}\label{glyph:approve}
Approve.  To be right.
\end{glyph}
\gindex{Approve}{是}\strindex{9}{是}

\begin{comps}
\comprtwocone & 日 + 疋⺪ \\\hline
\comprthreecone & Sun/Day + Roll/Bolt of Cloth \\\hline
\end{comps}

\begin{remglyph}
\hooglig{Approve} \remcom{today} whether the \remcom{roll} of scripture is in order.\\\hline
\end{remglyph}

\begin{prons}
\pronsrtwocone & \textbf{ゼ (ZE)} & --- \\\hline
\pronsrthreecone & --- & --- \\\hline
\end{prons}

\begin{remprons}
\hooglig{Approve} of \comon{zealous} works.\\\hline
\end{remprons}

%{\middel}
% &  \mbox{()}\\\hline
\begin{somme}
是非 & approve + un- = \textit{by any means; without fail} & ゼ{\middel}ヒ \mbox{(ZE{\middel}HI)}\\\hline
\end{somme}

\begin{decomp}{4}
是非&彼&に&会いたい。\\\hline
ゼヒ&かれ&に&あいたい\\\hline
by any means&he&に&to meet\\\hline
\multicolumn{4}{|r|}{\textit{I want to see him at all costs.}}\\\hline
\end{decomp}

\end{multicols}
%%--------------------------------------------------------------------
%22-------------------------------------------------------------------
\begin{multicols}{2}
\begin{glyph}{Straight (直)}{straight}\label{glyph:straight}
Straight, direct, at once.\\\hline
By extension comes the idea of things being correct or corrected (best illustrated by the transitive/intransitive verb pair below).  An important aspect of these two verbs can be seen when they attach to the end of other 漢字 as a suffix (as in the third example).  In such cases, this character gives the compound a sense of `redoing' whatever the first 漢字 represents.\\\hline
Ensure that you can clearly see the difference between this 漢字 and 真 (``true'', from Entry \ref{glyph:true}).
\end{glyph}
\gindex{Straight}{直}\strindex{8}{直}

\begin{comps}
\comprtwocone & 丨 + 目 + 十 \\\hline
\comprthreecone & (Pole) Shelf + Eye + Ten \\\hline
\end{comps}

\begin{remglyph}
\hooglig{The straight} \remcom{pole} \remcom{impaled} \remcom{eye} of the cross \remcom{scarecrow}.\\\hline
\end{remglyph}

\begin{prons}
\pronsrtwocone & \textbf{チョク (CHOKU)} & \textbf{なお (nao)}\\\hline
\pronsrthreecone & ジキ (JIKI) & ただ (tada) \\\hline
\rye{3}{You will encounter the 訓読み in only one word: 直ちに (ただち　に/tadachi ni) \textit{immediately, at once}.}
\end{prons}

\begin{remprons}
\hooglig{The straight} \comon{chalk} from \ucomon{Tajikistan} was used by \comkun{Naomi} to write \ucomkun{tadah!}\\\hline
\end{remprons}

%{\middel}
% &  \mbox{()}\\\hline
\begin{somme}
直る (intr.) & \textit{to be corrected / repaired} & なお{\middel}る \mbox{(nao{\middel}ru)}\\\hline
治す (tr.) & \textit{to correct / repair / fix} & なお{\middel}す \mbox{(nao{\middel}su)}\\\hline
書き直す & write + straight = \textit{rewrite; write again} & か{\middel}き　なお{\middel}す \mbox{(ka{\middel}ki nao{\middel}su)}\\\hline
直線 & straight + line = \textit{straight line} & チョク{\middel}セン \mbox{(CHOKU{\middel}SEN)}\\\hline
\end{somme}

\begin{decomp}{5}
ここ&に&直線&を&引け。\\\hline
ここ&に&チョクセン&を&ひけ\\\hline
here&に&straight line&を&draw\\\hline
\multicolumn{5}{|r|}{\textit{Draw a straight line here.}}\\\hline
\end{decomp}

\end{multicols}
\vspace*{\fill}
\pagebreak
%%--------------------------------------------------------------------
%23-------------------------------------------------------------------
\vspace*{\fill}
\begin{multicols}{2}
\begin{glyph}{Illness (病)}{illness}\label{glyph:illness}
Illness.  Disease.
\end{glyph}
\gindex{Illness}{病}\strindex{10}{病}

\begin{comps}
\comprtwocone & 疒 + 丙 \\\hline
\comprthreecone & Sickness + Third (\ref{glyph:third}) \\\hline
\end{comps}

\begin{remglyph}
\hooglig{Illness}, \remcom{sickness}, and disease are \remcom{three} words meaning the same thing.\\\hline
\end{remglyph}

\begin{prons}
\pronsrtwocone & \textbf{ビョウ (BY\={O})} & --- \\\hline
\pronsrthreecone & ヘイ (HEI) & や、やまい (ya, yamai) \\\hline
\end{prons}

\begin{remprons}
\hooglig{Illness} makes \ucomkun{ya mind} go \ucomon{haywire} even when you're \ucomkun{young}.  Sooner or later it will \comon{be yours}.\\\hline
\end{remprons}

%{\middel}
% &  \mbox{()}\\\hline
\begin{somme}
\rye{3}{Note the final compounds below uses the less common 音読み for 持.}
病気 & illness + spirit = \textit{illness} & ビョウ{\middel}キ \mbox{(BY\={O}{\middel}KI)}\\\hline
病院 & illness + institution = \textit{hospital} & ビョウ{\middel}イン \mbox{(BY\={O}{\middel}IN)}\\\hline
持病 & hold + illness = \textit{chronic disease} & ジ{\middel}ビョウ \mbox{(JI{\middel}BY\={O})}\\\hline
\end{somme}

\begin{decomp}{5}
此の&病院&は&市立&です。\\\hline
この&ビョウイン&は&シリツ&です\\\hline
this&hospital&は&city&be / is\\\hline
\multicolumn{5}{|r|}{\textit{This hospital is run by the city.}}\\\hline
\end{decomp}

\end{multicols}
%%--------------------------------------------------------------------
%24-------------------------------------------------------------------
\begin{multicols}{2}
\begin{glyph}{Together (共)}{together}\label{glyph:together}
Together.
\end{glyph}
\gindex{Together}{共}\strindex{6}{共}

\begin{comps}
\comprtwocone & 艹 + 一　+ 八 \\\hline
\comprthreecone & Grass + One + Eight \\\hline
\end{comps}

\begin{remglyph}
\hooglig{Together}, the \remcom{grass} worshipped the \remcom{volcano} in \remcom{one} spirit.\\\hline
\end{remglyph}

\begin{prons}
\pronsrtwocone & \textbf{キョウ (KY\={O})} & \textbf{とも (tomo)}\\\hline
\pronsrthreecone & --- & --- \\\hline
\end{prons}

\begin{remprons}
\hooglig{Together} in \comon{Kyoto} \comkun{tomorrow}.\\\hline
\end{remprons}

%{\middel}
% &  \mbox{()}\\\hline
\begin{somme}
共 & \textit{together} & とも \mbox{(tomo)}\\\hline
共同 & together + same = \textit{cooperation; collaboration} & キョウ{\middel}ドウ \mbox{(KY\={O}{\middel}D\={O})}\\\hline
共和国 & together + harmony + country = \textit{republic} & キョウ{\middel}ワ{\middel}コク \mbox{(KY\={O}{\middel}WA{\middel}KOKU)}\\\hline
\end{somme}

\end{multicols}

\begin{decomp}{5}
アメリカ合衆国&は&共和国&で&ある。\\\hline
アメリカガッシュウコク&は&キョウワコク&で&ある\\\hline
United States&は&republic&で&ある\\\hline
\multicolumn{5}{|r|}{\textit{The United States is a republic.}}\\\hline
\end{decomp}

\vspace*{\fill}
\pagebreak
%%--------------------------------------------------------------------